\paragraph{IV 21, \S2: relative divisors on projective and multiprojective bundles (cote25, sites106--111).}
The section applies the classification of invertible Modules from1.1 to positive relative divisors. At sites106--107 the canonical direct-image formula marked$(*)$ converts the section defining a divisor into a map involving a symmetric power of $\underline E$; the author expresses fibrewise regularity through universal injectivity and the surjectivity of its transpose. Proposition2.2 describes the resulting correspondence with invertible quotients of the dual symmetric power, and Corollaries2.3--2.4 give the representing projective bundles and their sum over the degree. Paragraph2.5 introduces a finite family of projective bundles and their fibre product. Proposition2.6 and paragraph2.7 extend the classification and degree notation to this multiprojective bundle, explicitly treating the partial degrees as locally constant functions on the base. Sites108--109 connect this construction to the planned paragraph25 on the resultant and discriminant of forms. Paragraph2.8 invokes ChapterII, paragraph2, in the displayed direct-image formula$(xx)$; its cancelled alternative reference to Künneth is readable. Proposition2.9 and Corollary2.10 describe divisors of a fixed multidegree by invertible quotients of a tensor product of dual symmetric powers and give the associated representing projective bundle and sum. Paragraph2.11 then extracts the relative representability argument from the projective-bundle case and relates it to the Picard prescheme. At site110, under the displayed hypothesis on $f_*(\underline O_X)$, the proof uses ChapterIII, paragraph7, to introduce the finitely presented Module $\underline Q$ representing the relevant functor, identifies sections with homomorphisms to an invertible Module, and states the correspondence with sections of $\underline P(\underline Q)$. The proposition number on that page is typed2.11 with an uncertain pen correction resembling2 on its last digit; site111 explicitly refers to2.12. Theorem2.13 at site111 gives the open retrocompact representing sub-prescheme, and the text identifies the entire projective bundle in its geometrically integral-fibre case. Its proof invokes the openness of transverse regularity through the incomplete reference ``par.11 ...'' and obtains retrocompactness by reduction to the noetherian case. The transcription retains several visible features that affect reading: site106 says ``section de $\underline E$'' and uses $\underline O_X(\underline D)$; site107 prints ``décomposition de D'' and has a crowded extra subscript and crossed group in the symmetric-power expression of2.4; the displayed isomorphism of2.6 at site108 omits the $\underline M$ factor mentioned in the surrounding statement; and site109 prints $\underline O_X(D)$ and $f_*(\underline D)$ in the two formulas. These are preserved with the scanned alphabets and indices. The N/G variable at site110, the corrected proposition digit, and several typewriter overprints remain explicitly uncertain.
